\chapter{Literature Review}
\label{ch:literature-review}

\section{Conceptual Foundations}

This section should introduce the theories, definitions, and assumptions on which
the research depends. Organise the discussion by concepts or research questions
rather than presenting one summary paragraph per paper. Where definitions vary,
explain which interpretation the thesis adopts and why.

A strong review synthesises evidence across studies. It should identify areas of
agreement, explain conflicting findings, and distinguish limitations in available
data from limitations in methods or evaluation design.

\section{Representative Approaches}

Compare representative approaches using criteria that matter to the thesis, such
as data requirements, assumptions, reproducibility, external validation, and
computational cost. The comparison should make the rationale for the chosen
methodology transparent rather than merely declaring one family of methods to be
superior.

\begin{sidewaystable}[p]
  \centering
  \scriptsize
  \caption[Example wide literature comparison.]{Example wide literature comparison with placeholder data. A short, plain optional caption is supplied for the List of Tables, while this longer caption may contain the detail needed to interpret the table.}
  \label{tab:wide-review}
  \begin{tabularx}{0.92\textheight}{@{}L{31mm}*{7}{>{\centering\arraybackslash}X}@{}}
    \toprule
    \textbf{Study} & \textbf{Setting} & \textbf{Data} & \textbf{Design} & \textbf{Primary metric} & \textbf{External test} & \textbf{Code} & \textbf{Key limitation} \\
    \midrule
    Example A \parencite{example2024} & Laboratory & Public & Retrospective & Accuracy & No & Yes & Single-site evaluation \\
    Example B \parencite{example2025} & Simulation & Synthetic & Controlled & Error & Yes & No & Limited sensitivity analysis \\
    Example C & Field study & Mixed & Prospective & Utility & Partial & Yes & Small sample size \\
    Example D & Benchmark & Public & Cross-sectional & Score & No & Partial & Narrow task definition \\
    Proposed study & Multiple settings & Mixed & Reproducible & Several metrics & Yes & Yes & Replace with actual limitations \\
    \bottomrule
  \end{tabularx}
\end{sidewaystable}


\section{Research Gap}

Conclude the review by identifying a focused and defensible gap. State what remains
unknown, why existing evidence cannot resolve it, and how the proposed research
addresses it. Do not define the gap only as the absence of a particular technique;
connect it to a substantive research need.

As summarised in \Cref{tab:wide-review}, the research gap should follow from a critical synthesis rather than a catalogue of publications.
